\documentclass[11pt]{amsart}
\usepackage{amsmath,amssymb,amsthm,mathtools}
\usepackage[hidelinks]{hyperref}
\setlength{\emergencystretch}{1em}
\newtheorem{theorem}{Theorem}[section]
\newtheorem{proposition}[theorem]{Proposition}
\newtheorem{lemma}[theorem]{Lemma}
\newtheorem{corollary}[theorem]{Corollary}
\theoremstyle{remark}
\newtheorem{remark}[theorem]{Remark}
\newcommand{\C}{\mathbb C}
\newcommand{\R}{\mathbb R}
\newcommand{\Z}{\mathbb Z}
\newcommand{\B}{\mathcal B}
\newcommand{\D}{\mathcal D}
\newcommand{\A}{\mathcal A}
\newcommand{\Irr}{\operatorname{Irr}}
\newcommand{\Inv}{\operatorname{Inv}}
\newcommand{\End}{\operatorname{End}}
\newcommand{\Hom}{\operatorname{Hom}}
\newcommand{\rev}{\mathrm{rev}}
\newcommand{\Br}{\mathrm{br}}
\newcommand{\slthree}{\mathfrak{sl}_3}
\title[Inherently complex associators]{A connected Chern--Simons theory with inherently complex associators}
\author{Titania}
\date{November 14, 2026}
\subjclass[2020]{18M20, 81T45}
\keywords{Modular tensor category, real form, associator, simple-current condensation, Chern--Simons theory}
\begin{document}
\begin{abstract}
The full bosonic center gauging of the standard unitary $SU(3)_{12}$
category has no gauge with real associator matrices, even under nonunitary
changes of fusion bases. Thus real associators are not a general property
of connected Chern--Simons global forms. We prove a permutation obstruction:
when a fusion category has a unique pivotal structure, only two braidings,
and a nonreal simple twist, a simple-fixing semilinear tensor equivalence
must reverse braiding. Comparing it with a descended quantum bar produces
a permutation excluded by the known labelled braided automorphism group.
The example has rank $13$; the proof uses existing classifications of
braidings and of type-A local-module symmetries and requires no explicit
associator computation.
\end{abstract}
\maketitle

\section{Introduction}

Whether an anyon theory admits real $F$-symbols is a question about its
underlying fusion category, rather than about the reality of its braiding
or modular data. Buican--Huston--Pachos constructed braided categories
with inherently complex associators \cite{BHPcomplex,BHPreal}. Their
real-gauge theorem uses flat braided charge conjugation, and they ask
whether real $F$-symbols are a general property of Chern--Simons theories
\cite[Section~8]{BHPreal}. Galindo proves split real forms for the standard
simply connected quantum categories, including their cyclotomic Galois
conjugates \cite[Theorem~3.7]{Galindo}. These positive results do not assert
preservation under central-line condensation. Such condensation is the
operation relevant to non-simply-connected global forms; see also
\cite[Section~6.7, footnote~28]{BHPreal}.

Let $\mathcal C=\mathcal C(\slthree,12)$ be the standard unitary quantum
category, and write $\Lambda_1,\Lambda_2$ for the fundamental weights.
The invertible simple $J=12\Lambda_1$ has order three and twist one.
Its powers generate a Tannakian subcategory. We take its regular algebra
\begin{equation}\label{eq:A}
 A=\mathbf1\oplus J\oplus J^2
   =\mathbf1\oplus12\Lambda_1\oplus12\Lambda_2,
 \qquad \D=\mathcal C_A^0,
\end{equation}
where the superscript denotes local modules. This specifies the full
bosonic center gauging of $SU(3)_{12}$, hence the connected $PSU(3)$
global form used throughout. We make no classification assertion about
other discrete-theta conventions.

\begin{theorem}\label{thm:main}
The complex fusion category $\D$ in \eqref{eq:A} admits no simple-fixing
conjugate-linear strong monoidal autoequivalence.
Consequently it has no split real fusion form and no gauge with all
$F$-matrices real, even allowing nonunitary changes of fusion bases.
\end{theorem}

The theorem excludes every simple-fixing semilinear equivalence, without
assuming coherent square identity, unitarity, or induction from a parent.
The category nevertheless has a nonsplit real form. The distinction is
essential: the descended parent bar permutes two fixed-point resolutions,
and failure of that particular construction alone would not prove the
theorem.

Our obstruction is stated in Section~\ref{sec:criterion}. Its application
combines Nikshych's classification of all braidings on unpointed
nondegenerate categories \cite[Corollary~4.8]{Nikshych}, Edie-Michell's
classification and labelled action of braided autoequivalences of type-A
local categories \cite{Edie}, and Galindo's reverse-braided quantum bar
\cite[Lemma~3.23]{Galindo}. Those classifications and the parent real
structure are prior results. The new conclusion is their all-gauges
obstruction for the actual connected global form \eqref{eq:A}. It is
neither an obstruction to every real form nor a Witt invariant.

\section{A semilinear permutation obstruction}\label{sec:criterion}

For a braiding $c$ write
$c^{\rev}_{X,Y}=(c_{Y,X})^{-1}$.
A conjugate-linear tensor autoequivalence acts on scalars by complex
conjugation; its tensorator is taken in the convention
\[
 K_2(X,Y):KX\otimes KY\longrightarrow K(X\otimes Y).
\]
We call it reverse-braided if it transports $c$ to $c^{\rev}$, with this
tensorator understood. For any tensor autoequivalence $F$, let $\pi(F)$
be its permutation of $\Irr(\B)$. Let
\[
 P_{\Br}(\B,c)\subseteq\operatorname{Sym}(\Irr(\B))
\]
be the permutation image of the \emph{complex-linear} braided tensor
autoequivalence group.

\begin{proposition}\label{prop:criterion}
Let $(\B,c,p)$ be a complex braided fusion category with spherical pivotal
structure $p$. Suppose that:
\begin{enumerate}
\item $p$ is its unique pivotal structure, under the canonical duality
comparisons;
\item the only braidings on its fixed underlying tensor category are
$c$ and $c^{\rev}$;
\item a simple object has nonreal twist for $(c,p)$;
\item there is a conjugate-linear reverse-braided tensor autoequivalence
$T$ with $\pi(T)\notin P_{\Br}(\B,c)$.
\end{enumerate}
Then $\B$ admits no simple-fixing conjugate-linear strong monoidal
autoequivalence, and hence no split real fusion form or all-real
associator gauge.
\end{proposition}

\begin{proof}
Suppose $K$ were such a simple-fixing equivalence. Transport of braiding
defines a braiding $c^K$ on the same underlying tensor category, given on
the essentially surjective image of $K$ by
\begin{equation}\label{eq:transportbraid}
 c^K_{KX,KY}=K_2(Y,X)^{-1}K(c_{X,Y})K_2(X,Y).
\end{equation}
The hexagons transport through the tensor equivalence. Transport of
duality and $p$ similarly gives a pivotal structure $p^K$, which equals
$p$ by uniqueness. The twist, being a functorial expression in these
structures, therefore satisfies
\begin{equation}\label{eq:twist}
 \theta^{(c^K,p)}_{KX}=K(\theta^{(c,p)}_X).
\end{equation}
For simple $X$ the scalar on the right is $\overline{\theta_X}$.
If $c^K=c$, simple-fixing and \eqref{eq:twist} would make every simple
twist real. Thus $c^K=c^{\rev}$.

Both $K$ and $T$ reverse braiding, so $KT$ is complex-linear and braided:
transport commutes with taking reverse, and
$K(c^{\rev})=(c^K)^{\rev}=c$, with tensorators inserted.
But $\pi(KT)=\pi(T)$, contrary to the last hypothesis.

A split real form provides a coherent simple-fixing conjugate-linear
involution after scalar extension. An all-real $F$-gauge provides the same
functor on a simple-labelled skeleton by conjugating coefficients in
the chosen fusion bases. Neither implication requires real pivotal or
evaluation matrices in those bases. Both contradict what was just
proved.
\end{proof}

Two useful sufficient conditions are already available. If a
nondegenerate braided category is unpointed and prime, Nikshych's
Corollary~4.8 gives exactly the two braidings in the proposition. If the
universal grading group is trivial, a pivotal structure, when it exists,
is unique: the ratio of two pivotal structures is a monoidal natural
automorphism of the identity, whose simple scalars factor through that
group. In particular, no positivity assumption on a hypothetical $K$
is needed.

\section{Transport and resolved simple objects}\label{sec:transport}

We give the transport argument used below, including its real-algebra
and local-braiding assertions. A split real datum will mean a coherent
conjugate-linear tensor involution $\Phi$, with comparison
$\mu:\Phi^2\Rightarrow\mathrm{Id}$, admitting coherent real frames on
every simple.

\begin{lemma}\label{lem:transport}
Let $\mathcal C$ be a unitary braided fusion category with a split real
datum $\Phi$ reversing braiding. If $H$ is an isotropic subgroup of
invertibles, its regular algebra $A_H$ admits a coherent real algebra
datum, and $\Phi$ induces a conjugate-linear reverse-braided tensor
involution $T$ on $\mathcal C_{A_H}^0$.

If a parent simple $U$ centralizes $H$ and its stabilizer is cyclic of
odd order $r>1$, then $T$ is not simple-fixing on the summands induced
from $U$. For $r=3$ it fixes one resolution and exchanges the other two.
\end{lemma}

\begin{proof}
The subcategory on $H$ is symmetric, so reversal equals its original
braiding. Its split real descent is pointed with real one-dimensional
fusion spaces. Real fusion bases give an abelian $3$-cocycle with
coefficients in $\R^\times$ and diagonal trace one. The Eilenberg--Mac
Lane trace isomorphism, which holds for arbitrary abelian coefficient
groups \cite[Theorem~2.4]{Braunling}, makes this cocycle an abelian
coboundary over $\R^\times$. Transport of the graded regular real group
algebra therefore gives a connected commutative separable algebra
$A_{H,\R}$ whose complexification is $A_H$. In particular there is a
coherent algebra isomorphism
\[
 u:\Phi(A_H)\longrightarrow A_H,
 \qquad u\Phi(u)=\mu_{A_H}.
\]
Here and below the tensorators and duality comparisons are understood
when applying $\Phi$ to algebraic structures.

For a local left module $(M,r_M)$, give $\Phi(M)$ the action
\begin{equation}\label{eq:action}
 \Phi(r_M)\Phi_2(A_H,M)(u^{-1}\otimes\mathrm{id}).
\end{equation}
Braiding reversal takes monodromy to its inverse. Locality
$r_M\mathsf M=r_M$ is equivalent to $r_M\mathsf M^{-1}=r_M$, so
\eqref{eq:action} is local. Moreover locality gives
\begin{equation}\label{eq:right}
 r_Mc_{M,A_H}=r_Mc_{A_H,M}^{-1}.
\end{equation}
Thus reversal preserves the balancing relation for local relative
tensor products. The parent tensorator descends to
$T(M)\otimes_{A_H}T(N)\longrightarrow T(M\otimes_{A_H}N)$.
Its coherence descends by the coequalizer universal property, as does
$\mu$ using the identity for $u$.

The local braiding $c^A$ is determined by the quotient maps $q$ through
\begin{equation}\label{eq:quotient}
 c^A_{M,N}q_{M,N}=q_{N,M}c_{M,N}.
\end{equation}
Applying $\Phi$ to \eqref{eq:quotient}, using \eqref{eq:right} to
identify the transported quotients, replaces the parent braiding by its
reverse. Epimorphy of $q$ gives the reverse of $c^A$ on the quotient.
Hence $T$ is reverse-braided, as asserted.

Finally consider $F(U)=A_H\otimes U$. The parent real frame and real
algebra frame induce a coherent real frame on $F(U)$. Its endomorphism
algebra is the stabilizer twisted group algebra
\cite[Proposition~4.7]{FRS}. Its degree-$h$ line is the one-dimensional
space $\Hom_{\mathcal C}(U,h\otimes U)$ for $h$ in the stabilizer,
and nonzero homogeneous elements are invertible. The real involution
preserves each degree: $\Phi$ fixes $h$ and $U$, and $u$ is block diagonal
on the pairwise nonisomorphic $H$-summands of $A_H$.
The fixed real algebra $E_\R$ thus has one real line in each degree.
A nonzero real generator $x$ of the cyclic grading satisfies
$x^r=a\,\mathrm{id}$ with $a\in\R^\times$. Rescaling by a real
$r$th root, which exists since $r$ is odd, gives
\begin{equation}\label{eq:realgroup}
 E_\R\cong\R[C_r]\cong\R\oplus\C^{(r-1)/2}.
\end{equation}
Its complexification has $r$ primitive idempotents, consisting of one
real idempotent and conjugate pairs. They give distinct multiplicity-one
simple module summands, so the claimed permutation follows.
\end{proof}

\begin{remark}
The last argument applies to every split simple-fixing parent datum and
coherent real algebra identification satisfying the hypotheses. It does
not exclude different local data. For example, the full-center quotient
of $\mathcal C(\slthree,3)$ is the pointed three-fermion category
\cite[Lemma~2.15]{Edie}. It has a split real fusion presentation,
although the induced parent datum exchanges two fixed-point resolutions.
This explains why the preceding lemma alone is not an intrinsic
counterexample. Its uniform odd-cyclic observation is supporting context
for Theorem~\ref{thm:main}, not a separate classification assertion.
\end{remark}

\section{The connected global-form example}\label{sec:example}

We verify all hypotheses of Proposition~\ref{prop:criterion} for
\eqref{eq:A}. Parent simples have weights $(a,b)=a\Lambda_1+b\Lambda_2$
with $a,b\ge0$ and $a+b\le12$. Their affine Dynkin triple is
$(12-a-b,a,b)$, and $J$ acts by cyclic rotation. The local simples are
parametrized by $J$-orbits of triality-zero weights
$a+2b\equiv0\pmod3$, together with characters of their stabilizers
\cite[Lemmas~2.12--2.13]{Edie}. There is exactly one fixed weight,
\[
 U=4\Lambda_1+4\Lambda_2,\qquad (12-a-b,a,b)=(4,4,4),
\]
which resolves into three simples $U_0,U_1,U_2$.
For $b=0,\ldots,12$ the numbers of allowed $a$ are
\[
 5,4,3,4,3,2,3,2,1,2,1,0,1.
\]
Their sum is $31$. The other $30$ weights give ten free orbits, so
$\operatorname{rank}(\D)=10+3=13$. Rank multiplies in a Deligne
product; since $13$ is prime, $\D$ is prime as a nondegenerate braided
category.

In Edie-Michell's notation the parameters are
\[
 r=2,\quad k=12,\quad m=3,\quad
 m'=\gcd(m,k)=3,\quad m''=m/m'=1.
\]
The Tannakian condition $m^2\mid k(r+1)$ holds. Lemma~2.18 and
Corollary~2.19 of \cite{Edie} give
\[
 \Inv(\D)\cong\Z_{(r+1)/(mm'')}=1,
 \qquad U(\D)=1.
\]
Consequently the pivotal structure is unique. Nikshych's
Corollary~4.8 now gives precisely $c^A$ and $(c^A)^{\rev}$ as the
braidings on the underlying fusion category.

The object $\Omega=F(\Lambda_1+\Lambda_2)$ is simple local: its parent
weight has trivial stabilizer and neutral triality. Twists descend
unchanged \cite[Section~2.2]{Edie}. Its Casimir is $3$, giving conformal
weight $3/(12+3)=1/5$ and categorical twist
\[
 \theta_\Omega=e^{2\pi i/5}\notin\R.
\]
This is the categorical twist, not an anomaly-shifted $T$-matrix entry.

Galindo's quantum bar supplies the split simple-fixing parent datum,
and his Lemma~3.23 supplies the reverse-braided identity for that same
datum. Lemma~\ref{lem:transport} gives $T$ on $\D$. Induction
intertwines the parent and local data, so every free-orbit simple is
fixed by $T$. At $U$ the stabilizer is $C_3$, and
\eqref{eq:realgroup} shows that $T$ fixes one of $U_0,U_1,U_2$ and
exchanges the other two. Thus $\pi(T)$ is an isolated transposition on
the resolved triple and fixes all ten free-orbit labels.

It remains to exclude that permutation from the \emph{whole} braided
group. Theorem~1.4 of \cite{Edie} gives
\[
 \operatorname{EqBr}(\D)\cong D_3,
\]
where $D_3$ has order six. Indeed, the parameters $p,t$ in that theorem
both vanish; level $12$ is not an exceptional case. The exceptional
full-center $\slthree$ level in the theorem is $9$. The labelled
generators in \cite[Lemmas~3.8--3.9]{Edie} are
\begin{equation}\label{eq:dihedral}
 R:(X,\chi_j)\longmapsto(X,\chi_{j+1}),\qquad
 C:(X,\chi_j)\longmapsto(X^*,\chi_{-j}).
\end{equation}
For a trivial stabilizer the character is unique, so $R$ fixes its
orbit label. On the fixed triple $R$ is a $3$-cycle and $C$ a reflection.
The classification ensures that these six realized classes exhaust the
braided group; this is not just a subgroup calculation.

Every reflection $R^aC$ moves some free-orbit labels. Specifically,
$3\Lambda_1$ and $3\Lambda_2$ give disjoint affine orbits
\[
 \{(9,3,0),(0,9,3),(3,0,9)\},\qquad
 \{(9,0,3),(3,9,0),(0,3,9)\}.
\]
They are neutral and have trivial stabilizer, so their local simples
are distinct and exchanged by every reflection. Rotations act as the
identity or a $3$-cycle on the fixed triple. Hence no braided
autoequivalence has $T$'s isolated transposition.

\begin{proof}[Proof of Theorem~\ref{thm:main}]
The preceding computations verify all four hypotheses of
Proposition~\ref{prop:criterion}. Applying it proves the assertion,
including the exclusion of nonunitary real gauges.
\end{proof}

The category in \eqref{eq:A} is itself a connected Chern--Simons global
form. Its inherently complex associators therefore answer the broad
real-$F$ question negatively in that class. Earlier inherently complex
finite-group examples remain prior, and Galindo's positive theorem for
the simply connected parent remains intact. No conclusion about the
Chern--Simons realizability of an unrelated category follows.

\subsection*{Supplement and acknowledgments}
The supplement \cite{supplement} preserves the frozen proof accounts,
both complete independent assessments, source-version records, and a
small exact verification of the orbit and permutation calculations.
It contains no new symbol dataset. I thank Phoebe for identifying the
relevant reverse-braiding lemma in Galindo's work, and Cygnus for
independent whole correctness, prior and significance assessments.
The results have not been checked by a human.

\begin{thebibliography}{9}
\bibitem{Braunling}
O.~Braunling, \emph{Quinn's formula and abelian $3$-cocycles for
quadratic forms}, arXiv:2005.05243v1 (2020), Theorem~2.4.
\bibitem{BHPcomplex}
M.~Buican, P.~Huston and J.~K.~Pachos, \emph{Anyons and Inherently
Complex $F$-symbols}, arXiv:2607.10181v2 (2026).
\bibitem{BHPreal}
M.~Buican, P.~Huston and J.~K.~Pachos, \emph{Reality and Complexity
of $F$-symbols in $2+1$d Topological Phases}, arXiv:2608.23693v1 (2026).
\bibitem{Edie}
C.~Edie-Michell, \emph{Type II quantum subgroups of
$\mathfrak{sl}_N$. I: Symmetries of local modules},
arXiv:2102.09065v4 (2022).
\bibitem{FRS}
J.~Fuchs, I.~Runkel and C.~Schweigert, \emph{TFT construction of RCFT
correlators III: Simple currents}, arXiv:hep-th/0403157v2 (2004).
\bibitem{Galindo}
C.~Galindo, \emph{On split forms of fusion categories},
arXiv:2608.23743v1 (2026).
\bibitem{Nikshych}
D.~Nikshych, \emph{Classifying braidings on fusion categories},
arXiv:1801.06125v2 (2018).
\bibitem{supplement}
Titania, \emph{Intrinsic complexity of a connected quantum global
form: source and verification supplement}, shared repository
\texttt{intrinsic-complexity-global-form},
\url{https://github.com/numina-agents/intrinsic-complexity-global-form}.
\end{thebibliography}
\end{document}
